\documentclass[11pt,a4paper]{article}

\usepackage[margin=2.5cm]{geometry}
\usepackage{lmodern}
\usepackage[T1]{fontenc}
\usepackage{microtype}
\usepackage{amsmath,amssymb,mathtools}
\usepackage{booktabs}
\usepackage{array}
\usepackage{algorithm}
\usepackage{algpseudocode}
\usepackage{listings}
\usepackage{xcolor}
\usepackage[hidelinks]{hyperref}

\lstset{
  language=C++, basicstyle=\ttfamily\small, columns=fullflexible,
  keywordstyle=\color{blue!60!black}\bfseries, commentstyle=\color{gray}\itshape,
  frame=single, framerule=0.3pt, rulecolor=\color{gray!60}, xleftmargin=1em, xrightmargin=1em,
  showstringspaces=false, breaklines=true
}

\newcommand{\R}{\mathbb{R}}
\newcommand{\Aeq}{A_{\mathrm{eq}}}
\newcommand{\beq}{b_{\mathrm{eq}}}
\newcommand{\Ain}{A_{\mathrm{in}}}
\newcommand{\bin}{b_{\mathrm{in}}}
\newcommand{\norm}[1]{\lVert #1\rVert}
\newcommand{\pos}[1]{\left[#1\right]_+}
\DeclareMathOperator*{\argmin}{arg\,min}
\DeclareMathOperator*{\argmax}{arg\,max}

\title{Convex Optimization in \texttt{nabia}\\[0.3em]\large L-BFGS and the augmented Lagrangian method: theory, algorithms and implementation}
\author{nabia documentation}
\date{October 2026}

\begin{document}
\maketitle

\begin{abstract}
This document describes the algorithms in the \texttt{nabia} module
\texttt{optimization/convex}: the limited-memory BFGS method (\texttt{LBFGS})
for smooth unconstrained problems, and the augmented Lagrangian method, also called the
method of multipliers (\texttt{AugmentedLagrangian}), for smooth convex
problems with linear equality, inequality and box constraints. We derive the
augmented Lagrangian for equalities and inequalities and the multiplier updates. We then
discuss convergence in the convex case and the practical safeguards (penalty update,
inexact inner solves, stopping tests), and present L-BFGS as both a standalone
solver and the inner solver.
\end{abstract}

\tableofcontents

%==============================================================================
\section{Module overview}

\begin{table}[ht]
\centering
\begin{tabular}{@{}l>{\raggedright\arraybackslash}p{9.5cm}@{}}
\toprule
Type & Role \\
\midrule
\texttt{IOptimizer<T>} & Base interface: owns the parameters \texttt{T} through a \texttt{std::unique\_ptr} and declares \texttt{optimize(initial\_guess)}, which returns the solution. \\
\texttt{LBFGS} & \texttt{IOptimizer<LBFGSParams>}; unconstrained smooth minimization (Section~\ref{sec:lbfgs}). \\
\texttt{IConvexOptimization} & \texttt{IOptimizer<ConvexOptimizationParams>}; interface for constrained convex solvers. \\
\texttt{AugmentedLagrangian} & \texttt{IConvexOptimization}; linearly constrained convex minimization, using \texttt{LBFGS} for its subproblems. \\
\texttt{LinearConstraint} & $\Aeq x=\beq$ and $\Ain x\le\bin$ as Eigen matrices and vectors. \\
\texttt{BoxConstraints} & Bounds $\ell\le x\le u$ (Section~\ref{sec:box}). \\
\texttt{Constraints} & \texttt{std::variant<BoxConstraints, LinearConstraint>}; held as \texttt{std::optional} in \texttt{ConvexOptimizationParams}, where empty means unconstrained. \\
\bottomrule
\end{tabular}
\caption{Types in \texttt{nabia/optimization}.}
\end{table}

All solvers follow the same pattern: construct with a parameter object and the
objective (and gradient), then call
\texttt{optimize(std::optional<Eigen::VectorXd> initial\_guess)}. It returns
the solution and starts from zero when no guess is given. Diagnostics such as
multipliers, residuals and iteration counts are in \texttt{result()}. The same
solver can be run again from a different starting point. Invalid input (mismatched dimensions, inconsistent
constraints, nonpositive tolerances) raises \texttt{std::invalid\_argument}.

%==============================================================================
\section{Problem statement}

We consider
\begin{equation}\label{eq:problem}
  \min_{x\in\R^n} \; f(x)
  \quad\text{s.t.}\quad
  h(x) \coloneqq \Aeq x - \beq = 0, \qquad
  g(x) \coloneqq \Ain x - \bin \le 0,
\end{equation}
with $f:\R^n\to\R$ convex and continuously differentiable,
$\Aeq\in\R^{p\times n}$, $\Ain\in\R^{m\times n}$. Box constraints
$\ell \le x \le u$ are a special case of the inequalities
(Section~\ref{sec:box}).

\paragraph{Lagrangian and KKT conditions.}
With multipliers $\lambda\in\R^p$ for the equalities and $\mu\in\R^m$ for the
inequalities, the Lagrangian is
\[
  L(x,\lambda,\mu) = f(x) + \lambda^\top h(x) + \mu^\top g(x).
\]
Because the constraints are linear, no further constraint qualification is
needed: a point $x^\star$ solves \eqref{eq:problem} if and only if there exist
$(\lambda^\star,\mu^\star)$ satisfying the Karush--Kuhn--Tucker conditions
\begin{align}
  \nabla f(x^\star) + \Aeq^\top\lambda^\star + \Ain^\top\mu^\star &= 0
    && \text{(stationarity)} \label{eq:kkt-stat}\\
  \Aeq x^\star = \beq, \quad \Ain x^\star &\le \bin
    && \text{(primal feasibility)} \label{eq:kkt-feas}\\
  \mu^\star &\ge 0
    && \text{(dual feasibility)} \\
  \mu^\star_i\, g_i(x^\star) &= 0,\; i=1,\dots,m
    && \text{(complementarity)} \label{eq:kkt-comp}
\end{align}
The multipliers have a sensitivity interpretation: $\lambda^\star_i$ and
$\mu^\star_i$ are the rates of change of the optimal value with respect to the
right-hand sides $b_i$ (with opposite sign).

%==============================================================================
\section{From penalties to augmented Lagrangians}

\subsection{The quadratic penalty and its drawback}
The simplest way to remove the equality constraints is the quadratic penalty
\[
  Q_\rho(x) = f(x) + \frac{\rho}{2}\norm{h(x)}^2 .
\]
Its minimizer $x_\rho$ satisfies $\nabla f(x_\rho) + \Aeq^\top\big(\rho\,h(x_\rho)\big) = 0$.
Comparing with \eqref{eq:kkt-stat}, $\rho\,h(x_\rho)$ plays the role of the
multiplier. Since $\lambda^\star$ is generally nonzero, $h(x_\rho)\approx\lambda^\star/\rho$
vanishes only as $\rho\to\infty$. The Hessian of $Q_\rho$ contains the term
$\rho\,\Aeq^\top\Aeq$, so its condition number grows like $\rho$ and the
subproblems become increasingly hard to solve accurately.

\subsection{Shifting the penalty: equality constraints}
The augmented Lagrangian adds the penalty to the Lagrangian instead of to $f$:
\begin{equation}\label{eq:al-eq}
  L_\rho(x,\lambda) = f(x) + \lambda^\top h(x) + \frac{\rho}{2}\norm{h(x)}^2 ,
  \qquad
  \nabla_x L_\rho = \nabla f(x) + \Aeq^\top\big(\lambda + \rho\,h(x)\big).
\end{equation}
If $x^+$ minimizes $L_\rho(\cdot,\lambda)$, then
$\nabla f(x^+) + \Aeq^\top(\lambda+\rho h(x^+)) = 0$. This is exactly the
stationarity condition \eqref{eq:kkt-stat} with the multiplier estimate
\begin{equation}\label{eq:update-eq}
  \lambda^+ = \lambda + \rho\, h(x^+).
\end{equation}
Now $h(x^+) = (\lambda^+-\lambda)/\rho$, so feasibility only requires the
multiplier estimates to converge, not $\rho\to\infty$. If $\lambda=\lambda^\star$,
then $x^\star$ already minimizes $L_\rho(\cdot,\lambda^\star)$ for every
$\rho>0$: the penalty is \emph{exact} at the right multiplier.

\subsection{Inequality constraints}
Write each inequality with a slack, $g_i(x) + s_i = 0$, $s_i\ge 0$, and apply
\eqref{eq:al-eq} to the equality $g(x)+s=0$:
\[
  \min_{x,\; s\ge 0}\; f(x) + \mu^\top\big(g(x)+s\big) + \frac{\rho}{2}\norm{g(x)+s}^2 .
\]
The minimization over $s$ separates by component and has a closed form. For
fixed $x$, the unconstrained minimizer of
$\mu_i(g_i+s_i)+\tfrac{\rho}{2}(g_i+s_i)^2$ is $s_i=-g_i-\mu_i/\rho$; projecting
onto $s_i\ge 0$ gives $s_i^\star = \max\{0,\,-g_i-\mu_i/\rho\}$. Substituting back:
\begin{itemize}
  \item if $\mu_i+\rho g_i \le 0$ (constraint comfortably inactive), then
        $g_i+s_i^\star=-\mu_i/\rho$ and the term equals $-\mu_i^2/(2\rho)$;
  \item if $\mu_i+\rho g_i > 0$, then $s_i^\star=0$ and the term equals
        $\mu_i g_i+\tfrac{\rho}{2}g_i^2 = \big((\mu_i+\rho g_i)^2-\mu_i^2\big)/(2\rho)$.
\end{itemize}
Both cases combine into the Powell--Hestenes--Rockafellar (PHR) augmented
Lagrangian
\begin{equation}\label{eq:al}
  \boxed{\;
  L_\rho(x,\lambda,\mu) = f(x) + \lambda^\top h(x) + \frac{\rho}{2}\norm{h(x)}^2
  + \frac{1}{2\rho}\sum_{i=1}^m\Big(\pos{\mu_i+\rho\,g_i(x)}^2 - \mu_i^2\Big)
  \;}
\end{equation}
with $\pos{t}=\max\{0,t\}$, and gradient
\begin{equation}\label{eq:al-grad}
  \nabla_x L_\rho(x,\lambda,\mu)
  = \nabla f(x) + \Aeq^\top\big(\lambda+\rho\,h(x)\big) + \Ain^\top\pos{\mu+\rho\,g(x)} .
\end{equation}
Applying \eqref{eq:update-eq} to $g+s^\star$ gives the inequality update
\begin{equation}\label{eq:update-in}
  \mu^+ = \mu + \rho\big(g(x^+)+s^\star\big) = \pos{\mu + \rho\, g(x^+)} ,
\end{equation}
which keeps $\mu^+\ge 0$ automatically.

Properties of \eqref{eq:al}:
\begin{itemize}
  \item $L_\rho(\cdot,\lambda,\mu)$ is convex whenever $f$ is (the extra terms
        are convex functions of affine maps), and continuously differentiable.
  \item It is \emph{not} twice differentiable where $\mu_i+\rho g_i(x)=0$.
        This is why a quasi-Newton inner solver (L-BFGS) is a good fit: it
        only uses gradients.
  \item With the updated multipliers, the gradient
        \eqref{eq:al-grad} at $x^+$ is exactly the stationarity residual of the
        KKT system:
        $\nabla_x L_\rho(x^+,\lambda,\mu) = \nabla f(x^+)+\Aeq^\top\lambda^+ +\Ain^\top\mu^+$.
        The inner solver's gradient norm therefore measures KKT stationarity
        directly.
\end{itemize}

%==============================================================================
\section{Convergence in the convex case}

\paragraph{Dual interpretation.}
Let $d(\lambda,\mu)=\inf_x L(x,\lambda,\mu)$ be the (concave) dual function.
Rockafellar~\cite{rockafellar1976} showed that one step of the method of
multipliers, minimizing \eqref{eq:al} and applying
\eqref{eq:update-eq}--\eqref{eq:update-in}, is exactly one step of the
\emph{proximal point algorithm} applied to the dual:
\[
  (\lambda^+,\mu^+) = \argmax_{\lambda,\;\mu\ge 0}\;
    \Big\{ d(\lambda,\mu) - \frac{1}{2\rho}\norm{(\lambda,\mu)-(\lambda_k,\mu_k)}^2 \Big\}.
\]
Equivalently, it is a gradient ascent step of length $\rho$ on the smoothed dual,
whose gradient is the constraint residual. Consequences for convex problems:
\begin{itemize}
  \item For \emph{any} fixed $\rho>0$, if \eqref{eq:problem} has a KKT point, the
        multipliers converge to a dual solution and every limit point of
        $x_k$ is primal optimal. Unlike the pure penalty method, $\rho$ does not have to
        go to infinity.
  \item This remains true when the subproblems are solved inexactly, provided
        the errors are summable (e.g.\ inner tolerances $\varepsilon_k\to 0$
        geometrically).
  \item The rate is linear, and the linear constant improves as $\rho$ grows;
        with $\rho_k\to\infty$ it is superlinear. Larger $\rho$ makes outer
        iterations more effective but inner problems more ill-conditioned; the
        practical rule of Section~\ref{sec:practical} balances both.
\end{itemize}
For nonconvex $f$ the method still works locally (under second-order
sufficiency, for $\rho$ above a threshold), but \texttt{nabia} targets the
convex case.

%==============================================================================
\section{Practical algorithm}\label{sec:practical}

\paragraph{Measuring infeasibility.}
For the inequalities, $g_i\le 0$ alone is not enough: complementarity
\eqref{eq:kkt-comp} must hold too. We use the measure of Birgin and
Mart\'inez~\cite{birgin2014}
\begin{equation}\label{eq:violation}
  V_k = \max\Big\{ \norm{h(x_k)}_\infty,\;
                   \norm{\max\{g(x_k),\,-\mu_k/\rho_k\}}_\infty \Big\},
\end{equation}
evaluated \emph{before} the multiplier update. The $i$-th inequality term is zero
iff $g_i\le 0$ and either $g_i=0$ or $\mu_i=0$. Moreover,
$\max\{g,-\mu/\rho\} = (\mu^+-\mu)/\rho$, so $V_k$ also measures how much the
multipliers still move.

\paragraph{Penalty update.}
If the violation did not decrease enough,
$V_k > \tau\, V_{k-1}$ (default $\tau=0.25$), the multipliers alone are not
making fast enough progress and the penalty is increased,
$\rho_{k+1}=\min\{\gamma\rho_k,\rho_{\max}\}$ (default $\gamma=10$).
Otherwise $\rho$ is kept.

\paragraph{Inexact inner solves.}
Early outer iterations use inaccurate multipliers, so solving their subproblems
to high accuracy is wasted work. The inner tolerance starts at
$\varepsilon_0 = \max\{\texttt{tol},\,10^{-1}\}$ and shrinks as
$\varepsilon_{k+1}=\max\{\texttt{tol},\,0.1\,\varepsilon_k\}$. Each subproblem is warm
started at the previous $x_k$.

\paragraph{Stopping test.}
The method stops when both KKT residuals are below the user tolerance:
\[
  \norm{\nabla_x L_{\rho_k}(x_{k+1},\lambda_k,\mu_k)}_\infty \le \texttt{tol}
  \quad\text{and}\quad V_k \le \texttt{tol}.
\]
By the identity at the end of Section~2, the first condition is stationarity
\eqref{eq:kkt-stat} for the updated multipliers. The multipliers are
nonnegative by construction, and $V_k$ covers feasibility and complementarity.
So the test certifies an approximate KKT point.

\begin{algorithm}[ht]
\caption{Augmented Lagrangian (PHR) with L-BFGS inner solver}\label{alg:al}
\begin{algorithmic}[1]
\Require $x_0$, tolerance $\texttt{tol}$, $\rho_0>0$, $\gamma>1$, $\tau\in(0,1)$, $\rho_{\max}$
\State $\lambda_0\gets 0$, $\mu_0\gets 0$, $\varepsilon_0 \gets \max\{\texttt{tol},10^{-1}\}$, $V_{-1}\gets\infty$
\For{$k=0,1,\dots,\texttt{max\_iter}-1$}
  \State $x_{k+1}\gets\argmin_x L_{\rho_k}(x,\lambda_k,\mu_k)$ from $x_k$, up to $\norm{\nabla_x L_{\rho_k}}_\infty\le\varepsilon_k$ \Comment{Alg.~\ref{alg:lbfgs}}
  \State $V_k \gets \max\{\norm{h(x_{k+1})}_\infty,\ \norm{\max\{g(x_{k+1}),-\mu_k/\rho_k\}}_\infty\}$
  \State $\lambda_{k+1}\gets\lambda_k+\rho_k\,h(x_{k+1})$
  \State $\mu_{k+1}\gets\pos{\mu_k+\rho_k\,g(x_{k+1})}$
  \If{$\norm{\nabla_x L_{\rho_k}(x_{k+1},\lambda_k,\mu_k)}_\infty\le\texttt{tol}$ \textbf{and} $V_k\le\texttt{tol}$}
    \State \Return $(x_{k+1},\lambda_{k+1},\mu_{k+1})$ \Comment{approximate KKT point}
  \EndIf
  \If{$V_k > \tau V_{k-1}$}
    \State $\rho_{k+1}\gets\min\{\gamma\rho_k,\rho_{\max}\}$
  \Else
    \State $\rho_{k+1}\gets\rho_k$
  \EndIf
  \State $\varepsilon_{k+1}\gets\max\{\texttt{tol},\,0.1\,\varepsilon_k\}$
\EndFor
\end{algorithmic}
\end{algorithm}

When there are no constraints, \eqref{eq:al} reduces to $f$, and the
implementation solves a single subproblem directly to the tolerance
$\texttt{tol}$.

%==============================================================================
\section{The inner solver: L-BFGS}\label{sec:lbfgs}

Each subproblem $\min_x L_{\rho}(x,\lambda,\mu)$ is smooth and unconstrained.
\texttt{nabia} solves it with limited-memory BFGS~\cite{nocedal1980,liu1989},
which is also exposed as the standalone optimizer \texttt{LBFGS}.

\paragraph{Quasi-Newton direction.}
BFGS keeps an approximation $H_k$ of the inverse Hessian and moves along
$d_k=-H_k\nabla f(x_k)$. With $s_k=x_{k+1}-x_k$, $y_k=\nabla f(x_{k+1})-\nabla f(x_k)$
and $\rho_k=1/(y_k^\top s_k)$, the update is
\[
  H_{k+1} = (I-\rho_k s_k y_k^\top)\,H_k\,(I-\rho_k y_k s_k^\top) + \rho_k s_k s_k^\top .
\]
L-BFGS never forms $H_k$. It stores the last $M$ pairs $(s_j,y_j)$ and computes
$H_k\nabla f$ with the two-loop recursion (Algorithm~\ref{alg:lbfgs}) in
$O(Mn)$ time and memory. The initial matrix is
$H_k^0=\gamma_k I$ with $\gamma_k = s_{k-1}^\top y_{k-1}/y_{k-1}^\top y_{k-1}$,
which gives the step a meaningful scale.

\paragraph{Safeguards.}
\begin{itemize}
  \item A pair is stored only if $s^\top y > 10^{-12}\norm{s}\norm{y}$.
        This keeps $H_k$ positive definite and $d_k$ a descent direction.
        For convex $f$ the curvature condition $s^\top y\ge 0$ always holds.
  \item If $d_k$ is not a descent direction (for example due to round-off), the memory is
        cleared and a normalized steepest-descent step is taken.
  \item The step length $t$ is chosen by backtracking ($t\gets 0.5t$) from $t=1$
        until the Armijo condition
        $f(x_k+t d_k)\le f(x_k)+c_1 t\,\nabla f(x_k)^\top d_k$ holds
        (default $c_1=10^{-4}$). Non-finite trial values are rejected.
\end{itemize}

\begin{algorithm}[ht]
\caption{L-BFGS with Armijo backtracking}\label{alg:lbfgs}
\begin{algorithmic}[1]
\Require $x_0$, memory $M$, tolerance $\varepsilon$
\For{$k=0,1,\dots$}
  \If{$\norm{\nabla f(x_k)}_\infty\le\varepsilon$} \Return $x_k$ \EndIf
  \State $q\gets\nabla f(x_k)$
  \For{$j=k-1,\dots,k-M$} \Comment{newest to oldest}
    \State $\alpha_j\gets\rho_j s_j^\top q$; \quad $q\gets q-\alpha_j y_j$
  \EndFor
  \State $r\gets\gamma_k q$
  \For{$j=k-M,\dots,k-1$} \Comment{oldest to newest}
    \State $\beta\gets\rho_j y_j^\top r$; \quad $r\gets r+(\alpha_j-\beta)s_j$
  \EndFor
  \State $d_k\gets -r$; \textbf{if} $\nabla f(x_k)^\top d_k\ge 0$ \textbf{then} clear memory, $d_k\gets-\nabla f(x_k)/\max\{1,\norm{\nabla f(x_k)}\}$
  \State $t\gets 1$; \textbf{while} $f(x_k+td_k)>f(x_k)+c_1t\nabla f(x_k)^\top d_k$ \textbf{do} $t\gets t/2$
  \State $x_{k+1}\gets x_k+td_k$; store $(s_k,y_k)$ if $s_k^\top y_k>0$, dropping the oldest pair beyond $M$
\EndFor
\end{algorithmic}
\end{algorithm}

%==============================================================================
\section{Box constraints}\label{sec:box}

\texttt{BoxConstraints} $\ell\le x\le u$ are turned into rows of $\Ain x\le\bin$.
Every finite $u_i$ gives the row $e_i^\top x\le u_i$, and every finite $\ell_i$ gives
$-e_i^\top x\le-\ell_i$. Infinite bounds produce no row, and an empty bound
vector means that side is unbounded. In \texttt{result().mu\_ineq} the upper-bound
multipliers come first (in index order), followed by the lower-bound
multipliers.
Bounds with $\ell_i>u_i$ are rejected with \texttt{std::invalid\_argument}.

%==============================================================================
\section{Using the implementation}

\subsection{Augmented Lagrangian}
\begin{lstlisting}
#include <nabia/optimization/convex/AugmentedLagrangian.h>

// min |x - c|^2  s.t.  sum(x) = 1, x >= 0   (projection onto the simplex)
Eigen::VectorXd c(3);  c << 1.0, -1.0, 0.5;

LinearConstraint lc;
lc.A_eq   = Eigen::MatrixXd::Ones(1, 3);      lc.b_eq   = Eigen::VectorXd::Ones(1);
lc.A_ineq = -Eigen::MatrixXd::Identity(3, 3); lc.b_ineq = Eigen::VectorXd::Zero(3);

auto params = std::make_unique<ConvexOptimizationParams>(
    ConvexOptimizationParams{3, /*tolerance*/ 1e-8, /*max_iter*/ 100, lc});

AugmentedLagrangian solver(
    std::move(params),
    [&](const Eigen::VectorXd& x) { return (x - c).squaredNorm(); },
    [&](const Eigen::VectorXd& x) -> Eigen::VectorXd { return 2.0 * (x - c); });

Eigen::VectorXd x = solver.optimize(std::nullopt);  // x = (0.75, 0, 0.25)
const auto& r = solver.result();   // r.lambda_eq = (0.5), r.mu_ineq = (0, 2.5, 0)
\end{lstlisting}
The KKT conditions can be checked by hand:
$2(x-c)+\lambda\mathbf{1}-\mu = 0$ gives $\lambda=0.5$ from the first
component (where $\mu_1=0$) and $\mu_2 = 2(0+1)+0.5=2.5$ for the active bound
$x_2=0$.

\begin{table}[ht]
\centering
\begin{tabular}{@{}llp{6.5cm}@{}}
\toprule
Option & Default & Meaning \\
\midrule
\texttt{initial\_penalty} & $10$ & $\rho_0$ \\
\texttt{penalty\_growth} & $10$ & $\gamma$, factor applied to $\rho$ when progress stalls \\
\texttt{max\_penalty} & $10^{10}$ & $\rho_{\max}$ \\
\texttt{required\_decrease} & $0.25$ & $\tau$ in the test $V_k>\tau V_{k-1}$ \\
\texttt{initial\_inner\_tolerance} & $10^{-1}$ & $\varepsilon_0$ (never below \texttt{tolerance}) \\
\texttt{max\_inner\_iter} & $1000$ & L-BFGS iterations per subproblem \\
\texttt{lbfgs\_memory} & $10$ & $M$, number of stored pairs \\
\bottomrule
\end{tabular}
\caption{Fields of \texttt{AugmentedLagrangianOptions}.}
\end{table}

The result type \texttt{AugmentedLagrangianResult} reports \texttt{x},
\texttt{lambda\_eq}, \texttt{mu\_ineq}, \texttt{objective},
\texttt{constraint\_violation} ($V_k$),
\texttt{stationarity}, the final \texttt{penalty}, the outer and total inner
iteration counts, and \texttt{converged}.

\subsection{L-BFGS on its own}
\begin{lstlisting}
#include <nabia/optimization/convex/LBFGS.h>

// Rosenbrock, using a single callback for value and gradient.
auto fg = [](const Eigen::VectorXd& z, Eigen::VectorXd& grad) {
    const double x = z(0), y = z(1);
    grad(0) = -2.0 * (1.0 - x) - 400.0 * x * (y - x * x);
    grad(1) = 200.0 * (y - x * x);
    return (1.0 - x) * (1.0 - x) + 100.0 * (y - x * x) * (y - x * x);
};
Eigen::VectorXd x0(2);  x0 << -1.2, 1.0;

LBFGS solver(std::make_unique<LBFGSParams>(LBFGSParams{2}), fg);
Eigen::VectorXd x = solver.optimize(x0);           // x = (1, 1)
if (!solver.result().converged()) { /* inspect solver.result().status */ }
\end{lstlisting}
\texttt{LBFGSParams} holds \texttt{dim}, \texttt{tolerance} (default
$10^{-8}$), \texttt{max\_iter} ($1000$), \texttt{memory} ($10$) and the line-search
constants \texttt{armijo\_c1}, \texttt{backtrack}, \texttt{max\_line\_search}.
\texttt{LBFGSResult::status} is \texttt{Converged}, \texttt{MaxIterations} or
\texttt{LineSearchFailed}.

\subsection{Scaling advice}
The method is invariant to neither the scaling of $f$ nor that of the
constraint rows. If the rows of $\Aeq,\Ain$ have very different norms, or $f$
takes values far from $O(1)$, normalize them first, since $\rho_0=10$ and
the tolerances assume a well-scaled problem. If the solver ends with
\texttt{converged == false} and a large \texttt{penalty}, the problem may be
infeasible or badly scaled.

%==============================================================================
\begin{thebibliography}{9}
\bibitem{hestenes1969}
M.~R. Hestenes, Multiplier and gradient methods,
\emph{J. Optim. Theory Appl.} 4 (1969) 303--320.

\bibitem{powell1969}
M.~J.~D. Powell, A method for nonlinear constraints in minimization problems,
in R.~Fletcher (ed.), \emph{Optimization}, Academic Press, 1969, 283--298.

\bibitem{rockafellar1973}
R.~T. Rockafellar, A dual approach to solving nonlinear programming problems by
unconstrained optimization, \emph{Math. Programming} 5 (1973) 354--373.

\bibitem{rockafellar1976}
R.~T. Rockafellar, Augmented Lagrangians and applications of the proximal point
algorithm in convex programming, \emph{Math. Oper. Res.} 1 (1976) 97--116.

\bibitem{bertsekas1982}
D.~P. Bertsekas, \emph{Constrained Optimization and Lagrange Multiplier Methods},
Academic Press, 1982.

\bibitem{nocedal1980}
J.~Nocedal, Updating quasi-Newton matrices with limited storage,
\emph{Math. Comp.} 35 (1980) 773--782.

\bibitem{liu1989}
D.~C. Liu and J.~Nocedal, On the limited memory BFGS method for large scale
optimization, \emph{Math. Programming} 45 (1989) 503--528.

\bibitem{nocedal2006}
J.~Nocedal and S.~J. Wright, \emph{Numerical Optimization}, 2nd ed., Springer,
2006, Chapters 7 and 17.

\bibitem{birgin2014}
E.~G. Birgin and J.~M. Mart\'inez, \emph{Practical Augmented Lagrangian Methods
for Constrained Optimization}, SIAM, 2014.
\end{thebibliography}

\end{document}
